% Generated draft: requires mathematical and source review.
\documentclass[11pt]{article}
\usepackage[margin=25mm]{geometry}
\usepackage{fontspec}
\setmainfont{Latin Modern Roman}
\usepackage{amsmath,amssymb,mathtools,unicode-math}
\setmathfont{Latin Modern Math}
\usepackage{enumitem,needspace,xurl,hyperref}
\setlength{\parindent}{0pt}
\setlength{\parskip}{5pt}
\setlength{\emergencystretch}{3em}
\newcommand{\N}{\mathbb{N}}
\newcommand{\Z}{\mathbb{Z}}
\newcommand{\Q}{\mathbb{Q}}
\newcommand{\R}{\mathbb{R}}
\newcommand{\C}{\mathbb{C}}
\newcommand{\F}{\mathbb{F}}
\newcommand{\A}{\mathbb{A}}
\newcommand{\PP}{\mathbb{P}}
\newcommand{\E}{\mathbb{E}}
\newcommand{\eps}{\varepsilon}
\newcommand{\dd}{\,\mathrm{d}}
\begin{document}
\section*{UnsolvedMath ID: 700128}
\textbf{Source problem code:} catalogue entry
\subsection{Definitions and mathematical statement}
Let \(k\) be a field with \(\operatorname{char}(k)=p>0\), and let \(\overline{k}\) be an algebraic closure. Under the classical smooth interpretation of ``linear algebraic group,'' consider smooth affine group schemes \(G\) of finite type over \(k\) that are commutative, connected, and unipotent. Here connectedness refers to the underlying Zariski topological space, and unipotence means that, for some integer \(n\geq 1\), the base change \(G_{\overline{k}}\) admits a closed immersion into
\[
U_n=\{(a_{ij})\in\operatorname{GL}_{n,\overline{k}}:a_{ii}=1\text{ for all }i,\ a_{ij}=0\text{ for }i>j\}.
\]
An isomorphism over \(k\) is an isomorphism of algebraic groups defined over \(k\).

\textbf{Classification problem (conditional classical formulation).} For each field \(k\) of positive characteristic, classify these groups up to isomorphism over \(k\): determine representatives such that every such group is isomorphic over \(k\) to exactly one representative.

\textbf{Unresolved conventions.} The source does not explicitly specify whether ``algebraic group'' requires smoothness or permits nonreduced group schemes, nor whether isomorphisms must be defined over the given field or only over its algebraic closure. The formulation above explicitly adopts the smooth, over-\(k\) interpretation and requires confirmation against the intended source conventions. The source imposes no perfectness or algebraic closedness hypothesis on \(k\); none is added here. Geometric unipotence as defined above does not impose that \(G\) be split over \(k\).

The source describes the characteristic-zero case as not difficult and attributes a classification up to isogeny to C. Chevalley. These are supplied contextual assertions, not independently verified results here; the requested classification is up to isomorphism.
\subsection{Short English statement}
For a field of positive characteristic, classify commutative connected unipotent linear algebraic groups up to isomorphism. A classical formulation using smooth groups and isomorphisms over the given field is stated above, but those conventions require confirmation. The supplied registry labels the problem open; this draft does not independently verify that status.
\subsection{Sources}
\begin{itemize}
\item Local UnsolvedMath registry: supplied complete source record, permanent ID \(700128\), exact source problem code \texttt{catalogue entry}, title \emph{Kourovka Notebook Problem 2.56}. The supplied status is \texttt{open}; the record was created and updated on \texttt{2026-10-16T14:00:00Z}, with upstream revision \texttt{d371ee917ac9647fbb06cd6fc529deec6f3dda82}. These are registry metadata, not evidence of independent mathematical or literature verification.
\item Supplied bibliographic attribution: \emph{The Kourovka Notebook (Unsolved Problems in Group Theory)}, problem 2.56, second issue (1967), as included in the twenty-first issue edition (2026), \texttt{arXiv:1401.0300v48}. Proposed by V. P. Platonov in 1967. The supplied collection description names E. I. Khukhro and V. D. Mazurov as editors.
\item The supplied statement attributes classification up to isogeny to C. Chevalley and describes classification in characteristic zero as not difficult. No separate reference or hypotheses for Chevalley's result are supplied. The cited Notebook and these contextual assertions have not been independently checked in preparing this draft.
\end{itemize}
\end{document}
